\documentclass[11pt]{article}

% -------------------------------------------------
% Packages
% -------------------------------------------------

\usepackage[a4paper, margin=1in]{geometry}
\usepackage{amsmath,amssymb}
\usepackage{graphicx}
\usepackage{booktabs}
\usepackage{hyperref}
\usepackage{caption}
\usepackage{subcaption}
\usepackage{siunitx}
\usepackage{setspace}

\setlength{\parindent}{0pt} % no indentation

% Citations
\usepackage[
backend=biber,
style=numeric-comp,
sorting=none
]{biblatex}

\addbibresource{references.bib}

% -------------------------------------------------
% Title
% -------------------------------------------------

\title{Project Title}

\author{
Student 1 \\
Student 2 \\
Student 3 \\
Student 4 \\
Student 5
}

\date{\today}

% -------------------------------------------------

\begin{document}

\maketitle

\begin{abstract}

Write a short summary (150-200 words) describing:
\begin{itemize}
\item The scientific question addressed in this project
\item The dataset used
\item The model or method applied
\item The main results
\end{itemize}

The abstract should allow a reader to understand the purpose and
outcome of the project without reading the full report.

\end{abstract}

\tableofcontents
\newpage

% -------------------------------------------------
\section{Introduction}

The introduction should explain the scientific context of the project.

Include the following elements:
\begin{itemize}
\item The broader research area (e.g., cosmology, particle physics, etc.)
\item Why the problem studied in this project is scientifically interesting
\item A brief overview of the physical phenomenon or experiment
\item What observable quantity is measured
\item The goal of this project
\end{itemize}

End the introduction with a short paragraph describing the structure of the report.

% -------------------------------------------------
\section{Physical Background}

Describe the physics relevant to the project.

Possible elements include:
\begin{itemize}
\item The theoretical framework
\item Important equations or physical relationships
\item The meaning of the observable quantities
\item The physical interpretation of the measurements
\end{itemize}

Explain the concepts clearly so that a reader familiar with
undergraduate physics can understand the analysis. Provide references to relevant literature where appropriate \cite{example}.


% -------------------------------------------------
\section{Data}

Describe the dataset used in the analysis.

Include:
\begin{itemize}
\item The source of the data
\item The format of the dataset
\item The measured quantities
\item Any preprocessing or filtering steps applied
\end{itemize}

If relevant, describe:
\begin{itemize}
\item uncertainties in the measurements
\item covariance between measurements
\item selection cuts applied to the data
\end{itemize}

You may include a summary table of the dataset.

% -------------------------------------------------
\section{Model}

Define the model used to describe the data.

Explain:
\begin{itemize}
\item the parameters of the model
\item the physical meaning of each parameter
\item the assumptions made
\end{itemize}

Write the model prediction as a function of parameters:
\begin{equation}
    m(\theta)
\end{equation}
where $\theta$ represents the vector of model parameters.

Clearly explain what features of the data the model aims to capture.

Also discuss any approximations or simplifications made in the model.

% -------------------------------------------------
\section{Statistical Method}

Explain the statistical framework used to analyze the data.

Describe how the model predictions are compared with the observed data.

If applicable, explain the assumptions made about measurement errors
(e.g., Gaussian uncertainties).

\subsection{Likelihood}

Define the likelihood function used in the analysis.

For example, a common choice is a Gaussian likelihood:
\begin{equation}
\ln \mathcal{L}(\theta)
=
-\frac{1}{2}(d-m(\theta))^{T}C^{-1}(d-m(\theta))
\end{equation}
where
\begin{itemize}
\item $d$ represents the data
\item $m(\theta)$ is the model prediction
\item $C$ is the covariance matrix.
\end{itemize}

Explain any simplifications made (for example using diagonal errors).

\subsection{Priors}

If Bayesian inference is used, specify the prior distributions
for the model parameters.

Explain:
\begin{itemize}
\item the choice of prior ranges
\item whether the priors are uniform, Gaussian, etc.
\item the physical motivation for these choices
\end{itemize}

% -------------------------------------------------
\section{Inference Procedure}

Describe how the posterior distribution of parameters is obtained.

Include details such as:
\begin{itemize}
\item the sampling algorithm used (e.g., MCMC)
\item the software library used
\item the number of samples generated
\item burn-in removal
\item convergence diagnostics
\end{itemize}

Explain how posterior samples are used to estimate parameter
uncertainties.

% -------------------------------------------------
\section{Results}

Present the main results of the analysis.

This section should include:
\begin{itemize}
\item posterior distributions of parameters
\item best-fit parameter values
\item credible intervals
\end{itemize}

\subsection{Parameter Constraints}

Include a summary table:
\begin{table}[h]
\centering
\begin{tabular}{ccc}
\toprule
Parameter & Estimate & Uncertainty \\
\midrule
Parameter 1 & ... & ... \\
Parameter 2 & ... & ... \\
\bottomrule
\end{tabular}
\caption{Posterior parameter estimates.}
\end{table}

\subsection{Model Comparison With Data}

Include a figure showing:
\begin{itemize}
\item the observed data
\item the model prediction
\item residuals if relevant
\end{itemize}

\begin{figure}[ht]
\centering
\includegraphics[width=0.8\textwidth]{example-image-a}
\caption{Comparison between data and model prediction.}
\end{figure}

Interpret the results and discuss the level of agreement.

% -------------------------------------------------
\section{Discussion}

Discuss the implications and limitations of the analysis.

Topics may include:
\begin{itemize}
\item sensitivity of results to modeling assumptions
\item statistical limitations
\item systematic uncertainties
\item comparison with results in the literature
\end{itemize}

% -------------------------------------------------
\section{Possible Extensions}

Describe possible improvements or extensions to the analysis.

Examples:
\begin{itemize}
\item using a more detailed model
\item incorporating additional datasets
\item testing sensitivity to prior choices
\item performing model comparison
\end{itemize}

% -------------------------------------------------
\section{Conclusion}

Summarize the main findings of the project.

Restate:
\begin{itemize}
\item the problem studied
\item the approach taken
\item the key results
\end{itemize}

Briefly comment on the significance of the results.

% -------------------------------------------------
\section*{Author Contributions}

Provide a short paragraph describing each student's contribution
to the project.

For example:
\begin{itemize}
\item Student A: data processing
\item Student B: model development
\item Student C: statistical inference
\item Student D: visualization
\item Student E: report writing and integration
\end{itemize}

% -------------------------------------------------

\printbibliography

% -------------------------------------------------
\appendix

\section{Additional Figures}

Include supplementary plots such as:
\begin{itemize}
\item histograms of the data
\item trace plots
\item diagnostic plots
\item additional comparisons
\end{itemize}

\section{Reproducibility}

Provide information needed to reproduce the analysis:
\begin{itemize}
\item software used
\item package versions
\item instructions for running the code
\item link to GitHub repository
\end{itemize}

\section{Feedback}
How doable was this project scale 1-10.

\end{document}